> Mía:
[4/20, 08:37] Kashana 🌻: % !TEX root = 00_instructivo.tex
\chapter{Transductores eléctricos y magnéticos: Magnetohidrodinamica}

\section*{Objetivos}
\begin{itemize}
    \item Visualizar campos magnéticos mediante limaduras de hierro entre dos placas acrílicas.
    \item Medir la intensidad del campo magnético generado por un imán y una PCB con corriente usando un magnetómetro QMC5883L.
    \item Cuantificar la corriente eléctrica en un circuito mediante el efecto Hall (sensor ACS712).
    \item Relacionar la corriente eléctrica, el campo magnético generado y su efecto sobre partículas ferrosas.
\end{itemize}

\section*{Preguntas previas}
\begin{enumerate}
    \item ¿Qué establece la ley de Ampere sobre la relación entre corriente eléctrica y campo magnético?
    \item ¿Cuál es la diferencia entre el efecto Hall y la magnetorresistencia anisotropa (AMR)?
    \item ¿Por qué es necesario calibrar un magnetómetro para compensar el \textit{hard-iron effect}?
    \item ¿Qué precauciones de seguridad se deben tener al trabajar con PCB y agua en el mismo experimento?
\end{enumerate}

\section*{Materiales y equipo}
\begin{itemize}
    \item Arduino Uno.
    \item Sensor de efecto Hall de corriente ACS712 (versión 5A).
    \item Magnetómetro de 3 ejes QMC5883L / GY-271.
    \item PCB con circuito integrado (pistas de cobre).
    \item Dos placas de acrílico transparente.
    \item Polvo ferroso (limaduras de hierro).
    \item Imanes de neodimio (diferentes tamaños).
    \item Fuente de alimentación externa para la PCB (batería 9V o fuente de laboratorio).
    \item Recipiente de vidrio con agua y tinte (opcional, para visualizar propulsión iónica).
    \item Mini protoboard y cables Dupont.
    \item Cinta aislante y pinzas de cocodrilo.
\end{itemize}

\section*{Procedimiento}
\begin{enumerate}
    \item \textbf{Preparación de la visualización magnética:}
    \begin{enumerate}
        \item Espolvoree una fina capa de polvo ferroso entre las dos láminas de acrílico. Selle los bordes con cinta.
        \item Coloque un imán debajo de las láminas y observe cómo se alinean las limaduras con las líneas de campo magnético.
        \item Dibuje o fotografíe el patrón observado.
    \end{enumerate}
    \item \textbf{Conexión y prueba del magnetómetro:}
    \begin{enumerate}
        \item Conecte el QMC5883L al Arduino (VCC a 3.3V, GND a GND, SDA a A4, SCL a A5).
        \item Cargue el código proporcionado en la guía (Sección 3.1).
        \item Acérquelo al imán y observe los valores en los ejes X, Y, Z. Anote la magnitud del campo.
    \end{enumerate}
    \item \textbf{Medición de corriente con el ACS712:}
    \begin{enumerate}
        \item Conecte el ACS712 (VCC a 5V, GND a GND, OUT a A1).
        \item Conecte la PCB (o un circuito simple con una resistencia) en serie con el terminal de medición del ACS712.
        \item Cargue el código de la Sección 3.2 y registre la corriente cuando la PCB está encendida y apagada.
    \end{enumerate}
    \item \textbf{Experimento integrador (propulsión iónica opcional):}
    \begin{enumerate}
        \item Coloque la PCB (con pistas desnudas) sumergida en el recipiente con agua y unas gotas de tinte.
        \item Conecte la PCB a la fuente externa a través del ACS712 para medir la corriente.
        \item Coloque el magnetómetro cerca de las placas.
        \item Encienda la corriente y observe la dirección del movimiento del tinte (fuerza de Lorentz).
        \item Registre la corriente medida y los cambios en el campo magnético.
    \end{enumerate}
\end{enumerate}

\section*{Esquemático}
\begin{figure}[htbp]
\centering
\begin{circuitikz}[scale=0.8, transform shape]
\draw (0,0) node[above]{Arduino Uno};
% Alimentación
\draw (1,0) node[ground]{GND};
\draw (1,0.5) node[anchor=south]{5V / 3.3V};
% Magnetómetro I2C

> Mía:
\draw (0.5,1) node[anchor=east]{A4 (SDA)} -- (2.5,1);
\draw (0.5,1.5) node[anchor=east]{A5 (SCL)} -- (2.5,1.5);
\node at (3,1.25) {QMC5883L (I2C)};
% ACS712
\draw (0.5,2.5) node[anchor=east]{A1} -- (2.5,2.5);
\node at (3,2.5) {ACS712 (OUT)};
% Circuito de potencia (PCB sumergida)
\draw (2.5,3.5) node[anchor=west]{PCB en agua} -- (1,3.5);
\draw (1,3.5) to[short, *-] (1,2.5);
\draw (1,2.5) to[short] (1,2.5) node[circ]{};
\draw (2.5,4) node[anchor=west]{Fuente externa} -- (1,4);
\draw (1,4) to[short] (1,3.5);
% El ACS712 va en serie con la PCB
\draw (1,2.5) node[anchor=east]{IP+} -- (1,3);
\draw (1,3) to[short] (1,3.5);
\end{circuitikz}
\caption{Conexiones para la práctica de transductores magnéticos. El ACS712 se conecta en serie con la carga (PCB) para medir la corriente que genera el campo magnético.}
\end{figure}

\section*{Preguntas de análisis}
\begin{enumerate}
    \item ¿Cómo se relaciona la magnitud del campo magnético medida por el QMC5883L con la corriente registrada por el ACS712?
    \item ¿Qué diferencias observa entre el patrón de limaduras generado por un imán y el generado por la PCB con corriente?
    \item ¿Por qué es importante usar una fuente externa y no los 5V del Arduino para alimentar la PCB sumergida?
    \item Explique físicamente por qué el tinte se mueve en una dirección específica al aplicar corriente en presencia de un campo magnético (Fuerza de Lorentz sobre los iones).
    \item ¿Qué fuentes de error pueden afectar las lecturas del magnetómetro y cómo se minimizan?
\end{enumerate}

\section*{Bibliografía}
\begin{itemize}
    \item Hoja de datos ACS712. Allegro MicroSystems.
    \item Hoja de datos QMC5883L. Honeywell.
   
\end{itemize}
[4/20, 08:37] Kashana 🌻: % !TEX root = 00_instructivo.tex
\chapter{Transductores mecánicos: Acelerómetro y podómetro}

\section*{Objetivos}
\begin{itemize}
    \item Comprender el principio de funcionamiento de un acelerómetro MEMS (MPU6050).
    \item Implementar un algoritmo de conteo de pasos (podómetro) usando la magnitud del vector de aceleración.
    \item Evaluar la precisión del podómetro en diferentes condiciones de movimiento (caminar, correr, movimientos bruscos).
    \item Alimentar el sistema con una batería portable para demostrar su uso en aplicaciones de monitoreo de actividad.
\end{itemize}

\section*{Preguntas previas}
\begin{enumerate}
    \item ¿Qué es una IMU (Unidad de Medición Inercial) y qué grados de libertad ofrece el MPU6050?
    \item ¿Por qué se usa la magnitud de la aceleración (\( \sqrt{a_x^2 + a_y^2 + a_z^2} \)) y no un solo eje para detectar pasos?
    \item ¿Qué es el \textit{efecto de histéresis} en el contexto de un umbral de detección?
    \item Investigue: ¿qué es el DMP (Digital Motion Processor) del MPU6050 y para qué sirve?
\end{enumerate}

\section*{Materiales y equipo}
\begin{itemize}
    \item Arduino Uno (o Arduino Nano para versión portable).
    \item Módulo MPU6050 (acelerómetro + giróscopo 6-DOF).
    \item Batería recargable de 9V o power bank 5V (para alimentación portable).
    \item Conector para batería (clips o cable USB).
    \item Mini protoboard.
    \item Cables Dupont.
    \item Cinta o velcro para fijar el conjunto a la cintura o al tobillo.
    \item Computadora con IDE Arduino (para carga inicial y monitoreo).
\end{itemize}

\section*{Procedimiento}
\begin{enumerate}
    \item \textbf{Conexión y prueba del MPU6050:}
    \begin{enumerate}
        \item Conecte el MPU6050 al Arduino (VCC a 5V, GND a GND, SDA a A4, SCL a A5).
        \item Cargue el código de la Sección 4.1 de la guía (podómetro).
        \item Verifique en el monitor serial que se leen valores de aceleración.
    \end{enumerate}
    \item \textbf{Calibración del umbral:}
    \begin{enumerate}
        \item Sujete el MPU6050 firmemente a su cintura o tobillo.
